\section{相关工作}
\label{sec:related-work-zh}

\subsection{异构图神经网络}

异构图神经网络通过建模节点类型、边类型和元路径语义来学习节点表示。
典型方法包括基于注意力的语义聚合、基于 Transformer 的异构消息传递，以及显式建模元路径实例的方法。

\subsection{鲁棒图神经网络}

鲁棒图神经网络研究关注图结构扰动对模型的影响，并通过边剪枝、图结构学习、概率建模或相似性约束提升防御能力。
这类方法大多面向同构图，直接迁移到异构图时可能忽略类型关系和多语义视图。

\subsection{鲁棒异构图学习}

鲁棒异构图学习进一步考虑异构结构中的关系类型和元路径语义。
已有方法通常在拓扑语义空间内识别可疑连接或调整语义权重。
DVCL 与这些方法的区别在于引入特征诱导视图，并通过跨视图对比学习降低对受攻击拓扑的单一依赖。

\subsection{图对比学习}

图对比学习通过构造不同视图或增强样本来学习更稳定的表示。
本文将对比学习用于异构图鲁棒防御场景，把净化拓扑视图和特征诱导视图作为互补视图进行对齐。
